%This is a template file for use of iopjournal.cls

\documentclass{iopjournal}

% Options
% 	[anonymous]	Provides output without author names, affiliations or acknowledgments to facilitate double-anonymous peer-review

\usepackage{amsmath}

\begin{document}

\articletype{Paper} %	 e.g. Paper, Letter, Topical Review...

\title{\textit{Physics-Informed Surrogate Model} Berbasis \textit{Partial}-LTE dan Persamaan Transfer Radiatif Dua-Zona untuk Diagnostik Plasma Bebas Kalibrasi}

\author{Birrul Walidain$^{1}$, Khairun Saddami$^2$,Rara Mitaphonna$^3$ and Nasrullah Idris$^{1,*}$}

\affil{$^1$Department of Physics, Faculty of Mathematics and Natural Sciences, Universitas Syiah Kuala, Banda Aceh, Indonesia}

\affil{$^2$Department of Electrical and Computer Engineering, Faculty of Engineering, Universitas Syiah Kuala, Banda Aceh, Indonesia}

\affil{$^3$Research Center for Photonics, National Research and Innovation Agency (BRIN), South Tangerang 15314, Indonesia}
% \affil{$^*$Author to whom any correspondence should be addressed.}

\email{nasrullah.idris@.usk.ac.id}

\keywords{LIBS, calibration-free, partial-LTE, two-zone model, radiative transfer, self-absorption, surrogate model}

\begin{abstract}
Analisis kuantitatif pada \textit{Calibration-Free Laser-Induced Breakdown Spectroscopy} (CF-LIBS) sangat bergantung pada asumsi plasma homogen berkesetimbangan termodinamika lokal (LTE). 
% Namun, plasma transien memiliki gradien spasial dan opasitas optik yang ekstrem, memicu penyerapan mandiri (\textit{self-absorption}) yang mendistorsi spektrum emisi. Lebih jauh, model kinetika kolisional-radiatif (CR) murni sering gagal akibat ketiadaan data penampang kolisional yang lengkap pada basis data standar (fenomena ``Gurun Data Atomik''). Untuk mengatasi keterbatasan ini, penelitian ini mengusulkan model maju (\textit{forward model}) dua-zona (Core-Shell) deterministik yang menggabungkan \textit{partial}-LTE (pLTE) dengan Persamaan Transfer Radiatif (RTE). Paradigma pLTE menstabilkan populasi eksitasi internal dari defisit data kolisional, sementara RTE secara analitik menyelesaikan efek \textit{self-reversal} makroskopis. Dengan mengintegrasikan parameter pelebaran Stark secara eksak dari literatur dibandingkan menggunakan korelasi probabilistik, model ini secara akurat mereproduksi depresi spektral menyerupai garis Fraunhofer pada transisi resonansi yang tebal secara optik. Arsitektur deterministik ini sepenuhnya menghindari algoritma \textit{machine learning} kotak-hitam pada mesin fisika dasar, menyediakan generator spektrum sintetis bebas kalibrasi yang ketat secara fisis. Kerangka kerja ini mereduksi ketidakpastian kuantitatif secara signifikan dalam menganalisis material bermatriks kompleks, menawarkan landasan yang tangguh untuk diagnostik plasma dinamis.

\end{abstract}

\section{Introduction}

Asumsi plasma homogen dalam Kesetimbangan Termodinamika Lokal (LTE) murni bertindak sebagai prasyarat fundamental bagi metode analisis kuantitatif \textit{calibration-free Laser-Induced Breakdown Spectroscopy} (CF-LIBS) \cite{favre_towards_2025, cristoforetti_local_2010}. Namun, prasyarat teoretis ini terbukti gagal untuk mendeskripsikan dinamika rekombinasi plasma riil yang bergradien tinggi \cite{favre_merlin_2025, zaitsev_two-zone_2024}. Dalam kondisi ideal, kerangka termodinamika ini memastikan parameter fundamental plasma, yakni temperatur elektron ($T_e$) dan densitas elektron ($n_e$), dapat diekstraksi secara presisi guna mengkuantifikasi komposisi unsur secara langsung, yang juga menjadi landasan utama dalam pemodelan campuran ionik pada plasma astrofisika \cite{fritzsche_atomic_2025}. Sebaliknya, penyimpangan dari kondisi LTE ini tidak hanya mendistorsi akurasi analitik, melainkan mereduksi mekanisme kuantifikasi menjadi tidak valid \cite{cristoforetti_local_2010, hansen_modeling_2021, zaytsev_accurate_2026}. Mengingat kompleksitas tersebut, perumusan karakterisasi plasma yang tervalidasi secara ketat menjadi sangat esensial, khususnya pada analisis material bermatriks kompleks yang secara inheren rentan terhadap anomali ketidakhomogenan spasial saat pengamatan, seperti yang diaplikasikan pada instrumen eksplorasi antariksa \cite{sawyers_database_2025, manelski_libs_2024}, analisis batuan pada kondisi atmosfer ekstrem \cite{zhang_improving_2025}, karakterisasi \textit{in-situ} transisi fase logam cair \cite{leosson_analysis_2022}, evaluasi abrasi pelindung panas wahana saat memasuki atmosfer \cite{leistikow_quantification_2026}, maupun pada pemantauan \textit{in-situ} dinding reaktor fusi nuklir \cite{marin_roldan_libs_2021}.

Upaya konvensional untuk menjustifikasi ekuilibrium absolut tersebut secara tradisional bersandar pada kriteria McWhirter (1965). Namun, perumusan klasik ini secara empiris terbukti menjadi syarat yang diperlukan tetapi tidak memadai untuk mendeskripsikan dinamika plasma riil \cite{bultel_mcwhirter_2025}. Kegagalan validasi ini berakar pada dua distorsi destruktif. Secara spasial, gradien termal memicu efek penyerapan mandiri (\textit{self-absorption}) yang merusak linearitas profil spektral emisi \cite{tang_review_2024}. Secara termodinamika, anomali populasi tingkat energi molekuler seperti pada intensitas transisi vibrasi-rotasi menunjukkan deviasi ekstrem dari prediksi LTE konvensional sehingga menuntut pendekatan multi-temperatur \cite{zaytsev_accurate_2026}. Untuk memitigasi anomali spasial ini, sejumlah studi komputasi mutakhir mulai mengintegrasikan arsitektur \textit{Artificial Neural Network} (ANN) dengan model dua-zona (\textit{two-zone model}) \cite{zaitsev_two-zone_2024, amogh_libs_2026,m_s_machine_2025}. Kendati pendekatan komputasi \textit{data-driven} ini berhasil mereduksi interferensi spasial, model-model tersebut mengalami kecacatan epistemologis karena masih tertahan pada asumsi kondisi stasioner. Padahal, kajian kinetika teoretis terbaru menegaskan bahwa plasma secara fundamental beroperasi secara sangat transien, di mana laju pendinginan yang masif dan tingginya opasitas radiatif ($\tau$) memaksa ambang batas densitas elektron melompat jauh melampaui prediksi stasioner \cite{bultel_mcwhirter_2025}. Kondisi ini mengekspos sebuah celah metodologis absolut. Arsitektur komputasi yang semata-mata mengandalkan resolusi spasial tanpa intervensi penyelesaian kinetik transien berisiko kuat memberikan validasi keliru (\textit{false validation}) pada spektrum yang sejatinya non-LTE. 

% Upaya komputasional termutakhir untuk merespons distorsi opasitas ini telah berevolusi, mulai dari penyelesaian analitik Persamaan Transfer Radiatif (RTE) \cite{favre_merlin_2025} hingga integrasi proksi \textit{machine learning} (ML) untuk mengakselerasi waktu komputasi \cite{favre_towards_2025}. Kendati inovasi metodologis ini berhasil mereduksi hambatan analitik secara signifikan, kerangka kerjanya secara umum masih bertumpu pada simplifikasi kondisi stasioner atau beroperasi sebagai ekstraktor emisi global tanpa resolusi spasial yang ketat. Keterbatasan dari asumsi stasioner ini semakin dipertegas oleh kajian teoretis terbaru yang menunjukkan bahwa perumusan ekuilibrium klasik—seperti kriteria McWhirter—cenderung mengunderestimasi ambang batas termodinamika pada plasma riil, terutama akibat tidak terakomodasinya laju pendinginan transien dan besarnya efek opasitas optik \cite{bultel_mcwhirter_2025}. Runtutan kompleksitas ini mengekspos sebuah celah metodologis yang esensial: pendekatan komputasi yang semata-mata mengandalkan pemecahan ruang spasial tanpa intervensi penyelesaian kinetik transien, maupun model prediktif yang mengabaikan kausalitas waktu, berisiko kuat memberikan validasi keliru terhadap spektrum plasma yang sejatinya berada pada rezim non-LTE.

% Upaya konvensional untuk menjustifikasi ekuilibrium absolut tersebut secara tradisional bersandar pada kriteria McWhirter (1965); namun, perumusan klasik ini secara empiris terbukti menjadi syarat yang diperlukan tetapi tidak memadai untuk mendeskripsikan dinamika plasma riil \cite{bultel_mcwhirter_2025}. Kegagalan validasi ini berakar pada dua distorsi destruktif: secara spasial, gradien termal memicu efek penyerapan mandiri (\textit{self-absorption}) yang merusak linearitas profil spektral emisi \cite{tang_review_2024}; dan secara termodinamika, deviasi populasi tingkat energi molekuler menunjukkan anomali ekstrem dari prediksi konvensional sehingga menuntut pendekatan multi-temperatur pada kondisi non-LTE \cite{terashkevich_two-temperature_2024}. Untuk merespons anomali spasial dan distorsi opasitas ini, instrumentasi komputasional termutakhir telah berevolusi, mulai dari penyelesaian analitik Persamaan Transfer Radiatif (RTE) \cite{favre_merlin_2025}, hingga pengintegrasian proksi \textit{machine learning} (ML) dan \textit{Artificial Neural Network} (ANN) berbasis model dua-zona guna mengakselerasi waktu komputasi \cite{zaitsev_two-zone_2024, amogh_libs_2026, favre_towards_2025}. Kendati inovasi metodologis ini berhasil mereduksi interferensi spasial secara signifikan, kerangka kerjanya secara umum masih memiliki keterbatasan komputasional mendasar karena tertahan pada simplifikasi kondisi stasioner atau beroperasi sebagai ekstraktor emisi global tanpa resolusi spasial yang ketat. Keterbatasan asumsi stasioner ini bertabrakan dengan realitas kinetika teoretis terbaru yang menegaskan bahwa plasma secara fundamental bersifat sangat transien; di mana laju pendinginan yang masif dan tingginya opasitas radiatif ($\tau$) memaksa ambang batas termodinamika melompat jauh melampaui prediksi stasioner \cite{bultel_mcwhirter_2025}. Runtutan kompleksitas ini mengekspos sebuah celah metodologis absolut: arsitektur komputasi yang semata-mata mengandalkan resolusi spasial tanpa intervensi penyelesaian kinetik transien, maupun model prediktif yang mengabaikan kausalitas waktu, berisiko kuat memberikan validasi keliru (\textit{false validation}) pada spektrum yang sejatinya non-LTE.


Menghadapi celah metodologis tersebut sekaligus merespons trilema antara kegagalan ekuilibrium analitik, kelumpuhan komputasi deterministik, dan kesesatan kausalitas \textit{machine learning}, penelitian ini merumuskan sebuah resolusi arsitektur \textit{Physics-Informed Surrogate Model}. Arsitektur ini menolak secara fundamental penggunaan algoritma prediksi sebagai pencocok pola buta, melainkan memformulasikannya untuk beroperasi secara eksklusif di atas ruang termodinamika yang telah dikendalikan (\textit{pre-constrained manifold}). Untuk mencapai hal tersebut, fase komputasi dibelah menjadi dua ranah absolut. Sebuah generator maju (\textit{forward model}) di fase \textit{offline} yang merajut termodinamika Dua-Zona \textit{partial}-LTE (pLTE) \cite{zaitsev_two-zone_2024} dengan penyelesaian analitik Persamaan Transfer Radiatif (RTE), serta sebuah proksi \textit{Support Vector Regression} (SVR) di fase \textit{online} \cite{wu_novel_2025}. Melalui pemisahan algoritmik yang deterministik ini, penelitian ini bertujuan untuk mengekstrak parameter suhu dan densitas plasma secara akurat dari spektrum emisi yang terdistorsi oleh \textit{self-absorption}, tanpa perlu menjalankan satu pun iterasi numerik RTE secara \textit{online}. Dalam integrasi hibrida ini, paradigma pLTE bertugas menstabilkan populasi eksitasi internal dari kegagalan ekuilibrium, sementara eksekusi SVR menjamin bahwa diagnostik tereksekusi secara seketika (\textit{real-time}) dengan tingkat kepatuhan absolut terhadap hukum fundamental mekanika kuantum. Sebuah formulasi yang secara spesifik dirancang untuk menaklukkan anomali fisis pada batas optik tebal (\textit{optically thick limit}) sekaligus mengeliminasi hambatan waktu komputasional pada spektroskopi analitik.


% Ekspansi plasma yang ekstrem memicu gradien spasial dan pendinginan cepat, yang secara simultan menghancurkan ekuilibrium termodinamika tingkat dasar secara fundamental \cite{zaytsev_accurate_2026} dan memicu efek opasitas optik berupa penyerapan mandiri (\textit{self-absorption}) \cite{hansen_modeling_2021}. Pendekatan simplifikasi analitik yang buta terhadap transfer radiatif ini akan secara keliru menginterpretasikan pelebaran profil spektral akibat opasitas sebagai murni efek Stark kolisional, yang bermuara pada overestimasi densitas elektron ($n_e$) dan depresiasi ekstraksi suhu ($T_e$). Upaya untuk mengoreksi anomali opasitas tersebut menghadirkan kebuntuan metodologis baru. Pendekatan inversi deterministik murni—yang secara langsung mengeksekusi iterasi Persamaan Transfer Radiatif (RTE)—memang berhasil mempertahankan integritas termodinamika secara absolut \cite{favre_merlin_2025}. Namun, penggunaan algoritma pencarian akar non-linier untuk membalikkan matriks RTE memicu hambatan komputasi (computational bottleneck) yang ekstrem, sehingga menggugurkan kelayakan metode analitik ini untuk keperluan diagnostik analitik \cite{favre_towards_2025}. Di sisi ekstrem lainnya, substitusi komputasi menggunakan algoritma machine learning konvensional menawarkan kecepatan, tetapi menyimpan cacat epistemologis. Sebagaimana dievaluasi oleh Hao et al. \cite{hao_machine_2024}, mayoritas algoritma murni data-driven beroperasi sebagai 'kotak-hitam' (black-box) stokastik yang mencederai kausalitas mekanika kuantum, memutus rantai hukum termodinamika dasar, dan memicu ketergantungan irasional pada akuisisi dataset empiris berskala masif.

% Menghadapi trilema antara kegagalan ekuilibrium analitik, kelumpuhan komputasi deterministik, dan kesesatan kausalitas \textit{machine learning}, penelitian ini merumuskan sebuah resolusi arsitektur \textit{Physics-Informed Surrogate Model}. Arsitektur ini menolak secara fundamental penggunaan \textit{machine learning} sebagai pencocok pola buta, melainkan memformulasikannya untuk beroperasi secara eksklusif di atas ruang termodinamika yang telah dikendalikan (\textit{pre-constrained manifold}). Untuk mencapai hal tersebut, fase komputasi dibelah menjadi dua ranah absolut: sebuah generator maju (\textit{forward model}) di fase \textit{offline} yang merajut termodinamika Dua-Zona \textit{partial}-LTE (pLTE) \cite{zaitsev_two-zone_2024} dengan penyelesaian analitik Persamaan Transfer Radiatif (RTE), serta sebuah proksi \textit{Support Vector Regression} (SVR) di fase \textit{online} \cite{wu_novel_2025}.

% Melalui pemisahan algoritmik yang deterministik ini, penelitian ini bertujuan untuk mengekstrak parameter suhu dan densitas plasma secara akurat dari spektrum emisi yang terdistorsi oleh \textit{self-absorption}, tanpa perlu menjalankan satu pun iterasi numerik RTE secara \textit{online}. Dalam arsitektur hibrida ini, paradigma pLTE bertugas menstabilkan populasi eksitasi internal dari kegagalan ekuilibrium, sementara penggunaan SVR menjamin bahwa ekstraksi diagnostik tereksekusi secara seketika (\textit{real-time}) dengan tingkat kepatuhan absolut terhadap hukum termodinamika fundamental. Integrasi ini secara spesifik dirancang untuk menaklukkan anomali fisis pada batas optik tebal (\textit{optically thick limit}), sekaligus mengeliminasi hambatan waktu komputasional yang selama ini melumpuhkan analisis spektroskopi analitik berbasis plasma.

% Asumsi plasma homogen yang selama ini digunakan dalam metode \textit{Calibration-Free Laser-Induced Breakdown Spectroscopy} (CF-LIBS) terbukti tidak memadai untuk mendeskripsikan dinamika plasma riil yang bergradien tinggi \cite{zaitsev_two-zone_2024}. Keandalan metode kuantitatif CF-LIBS mutlak bergantung pada pemenuhan syarat \textit{Local Thermal Equilibrium} (LTE) sebagai prasyarat utama \cite{cristoforetti_local_2010,favre_merlin_2025}. Dalam kondisi ideal, kerangka termodinamika ini memastikan parameter fundamental plasma, yakni temperatur elektron ($T_e$) dan densitas elektron ($n_e$), dapat diekstraksi secara presisi guna mengkuantifikasi komposisi unsur secara langsung, sebuah mekanisme fisis universal yang juga menjadi landasan utama dalam pemodelan campuran ionik pada plasma astrofisika \cite{fritzsche_atomic_2025}. Sebaliknya, penyimpangan dari kondisi LTE ini tidak hanya mendistorsi akurasi analitik, melainkan mereduksi mekanisme kuantifikasi menjadi tidak valid \cite{cristoforetti_local_2010,hansen_modeling_2021}. Mengingat kompleksitas pada karakterisasi plasma yang tervalidasi secara ketat menjadi sangat esensial, khususnya dalam analisis material matriks kompleks yang melekat rentan terhadap anomali ketidakhomogenan plasma pada pengamatan, seperti yang ditemui pada instrumen eksplorasi antariksa \cite{sawyers_database_2025,manelski_libs_2024}, analisis batuan pada kondisi atmosfer ekstrem \cite{zhang_improving_2025}, karakterisasi \textit{in-situ} transisi fase logam cair \cite{leosson_analysis_2022}, evaluasi abrasi pelindung panas wahana saat memasuki atmosfer \cite{leistikow_quantification_2026}, maupun pada pemantauan \textit{in-situ} dinding reaktor fusi nuklir \cite{marin_roldan_libs_2021}.

% Meskipun kriteria McWhirter (1965) secara tradisional digunakan sebagai landasan validasi LTE, perumusan klasik ini terbukti menjadi syarat yang diperlukan namun tidak memadai untuk mendeskripsikan dinamika plasma riil \cite{bultel_mcwhirter_2025}. Distorsi paling destruktif pada plasma adalah dalam bentuk penyerapan mandiri (\textit{self-absorption}), di mana struktur spasial yang berlapis merusak linearitas profil spektral emisi \cite{tang_review_2024}. Ketidakpastian ini semakin diperumit oleh anomali pada profil populasi tingkat energi internal molekuler di mana konfigurasi intensitas transisi vibrasi-rotasi menunjukkan deviasi signifikan dari prediksi LTE tradisional dan menuntut pendekatan diagnostik dua-temperatur pada kondisi non-LTE \cite{terashkevich_two-temperature_2024}. Untuk memitigasi anomali spasial tersebut, studi-studi komputasi terkini beralih pada arsitektur \textit{Artificial Neural Network} (ANN) yang mengintegrasikan model dua zona (\textit{two-zone model}) \cite{zaitsev_two-zone_2024,amogh_libs_2026,m_s_machine_2025}. Namun, sementara pendekatan komputasi \textit{data-driven} ini berhasil mereduksi interferensi spasial, model-model tersebut masih tertahan pada penyederhanaan kondisi stasioner. Padahal, kajian kinetika teoretis terbaru membuktikan bahwa plasma secara fundamental bersifat sangat transien; laju pendinginan yang ekstrem dan tingginya opasitas radiatif ($\tau$) memaksa ambang batas densitas elektron minimal melompat jauh melampaui prediksi stasioner konvensional \cite{bultel_mcwhirter_2025}. Kondisi ini menyisakan sebuah celah metodologis, di mana arsitektur analitik yang hanya mengandalkan pemecahan ruang spasial tanpa melibatkan filtrasi kinetik berbasis waktu berisiko kuat memvalidasi spektrum non-LTE secara keliru.

% Keterbatasan metode diagnostik konvensional seperti plot Saha-Boltzmann dalam menghadapi plasma transien telah diekspos secara empiris. Hansen \textit{et al.} \cite{hansen_modeling_2021} membuktikan bahwa asumsi plasma homogen gagal mengkompensasi distorsi opasitas, dan menegaskan bahwa penyederhanaan gradien spasial melalui arsitektur Transfer Radiatif (RTE) Dua-Zona adalah kerangka yang paling fisis dan akurat untuk memodelkan serapan mandiri (\textit{self-absorption}) tanpa batas hidrodinamika 3D. Lebih jauh, kegagalan ekuilibrium absolut di dalam plasma yang terekspansi dengan cepat juga telah dikonfirmasi oleh Zaitsev dan Terashkevich \cite{zaitsev_two-zone_2024}, yang menuntut terpisahnya derajat kebebasan fisis dan penerapan koreksi Kinetika Termodinamika (\textit{partial}-LTE). Oleh karena itu, penelitian ini merakit kedua fondasi fisik tersebut---RTE Dua-Zona dari Hansen sebagai penyelesai opasitas spasial dan Kinetika pLTE dari Zaitsev sebagai stabilisator eksitasi temporal---menjadi satu mesin generator deterministik. Pemisahan arsitektur ini melegitimasi penggunaan proksi interpolasi instan (\textit{Support Vector Regression}) di tahapan \textit{online}, selaras dengan pembuktian Liu \textit{et al.} \cite{lin_enhancing_2026} bahwa formulasi pencarian akar (\textit{root-finding}) berbasis komputasi heuristik atau optimasi multi-batasan ortodoks akan selalu lumpuh dan terjebak pada jebakan optimum lokal ketika dihadapkan pada redaman matriks inversi spektral.

% Upaya ekstraksi suhu ($T_e$) dan densitas elektron ($n_e$) dari plasma transien yang terdistorsi oleh opasitas optik menghadirkan sebuah trilema metodologis yang melumpuhkan paradigma konvensional. Ketika dihadapkan pada spektrum cacat akibat penyerapan mandiri (\textit{self-absorption}), peneliti umumnya dipaksa memilih satu dari tiga jalur komputasi yang tersedia, yang masing-masing menyimpan cacat fisis atau arsitektural yang fatal.

% Jalur pertama, yakni simplifikasi analitik yang diusung oleh model CF-LIBS tradisional, secara brutal melanggar realitas termodinamika. Pendekatan ini mengasumsikan transparansi optik absolut dan Kesetimbangan Termodinamika Lokal (LTE) murni \cite{cristoforetti_local_2010}. Akibatnya, pelebaran spektrum yang terjadi karena redaman opasitas secara keliru diinterpretasikan sebagai murni efek Stark kolisional, sehingga menghasilkan overestimasi densitas yang parah serta mengekstrak suhu dari tingkat energi yang sejatinya telah kehilangan ekuilibrium.

% Jalur kedua, yakni inversi deterministik murni, mempertahankan integritas termodinamika namun berujung pada kelumpuhan komputasional. Menggunakan algoritma pencarian akar non-linear secara iteratif untuk membalikkan Persamaan Transfer Radiatif (RTE) memicu hambatan komputasi (\textit{computational bottleneck}) yang ekstrem \cite{rodriguez_mixkiprapcal_2021,espinosa_analysis_2019}. Menjalankan integrasi matriks yang kaku pada data eksperimen beresolusi tinggi memakan waktu yang sangat tidak rasional, sehingga seketika menggugurkan potensi diagnostik secara \textit{real-time}.

% Jalur ketiga, yakni penerapan \textit{machine learning} konvensional (murni statistik), menawarkan kecepatan tinggi namun secara mutlak mencederai kausalitas mekanika kuantum \cite{m_s_machine_2025}. Algoritma ini terjebak dalam kesesatan kotak-hitam (\textit{black-box fallacy}), beroperasi murni sebagai pencocok pola tanpa pemahaman terhadap hukum termodinamika (Saha-Boltzmann). Ketergantungan absolut pada data eksperimen empiris membuat model ini rentan terhadap \textit{overfitting} dan melumpuhkan kemampuannya untuk melakukan generalisasi pada matriks material yang berbeda \cite{favre_towards_2025}.

% Pendekatan simplifikasi analitik pada model CF-LIBS konvensional sering kali mendistorsi realitas termodinamika dengan memaksakan parameter transparansi optik absolut dan Kesetimbangan Termodinamika Lokal (LTE) murni secara homogen \cite{cristoforetti_local_2010}. Ketika asumsi idealis ini dipertahankan pada plasma transien, kalkulasi analitik murni terbukti memicu deviasi kuantitatif yang parah terhadap realitas fisis \cite{hansen_modeling_2021}. Secara mekanistik, redaman foton akibat opasitas optik (self-absorption) menyebabkan profil spektral mengalami pelebaran artifisial sekaligus penurunan intensitas puncak. Akibatnya, sistem diagnostik yang buta terhadap transfer radiatif akan menyalahartikan pelebaran ini sebagai murni efek Stark kolisional, yang bermuara pada overestimasi densitas elektron $n_e$. Secara bersamaan, depresiasi intensitas emisi tersebut memaksa algoritma mengekstrak suhu dari tingkat energi yang sejatinya telah kehilangan ekuilibrium, sehingga menghasilkan underestimasi suhu analitik yang keliru. 

% Pendekatan inversi deterministik murni—yang secara langsung mengeksekusi iterasi Persamaan Transfer Radiatif (RTE)—memang berhasil mempertahankan integritas termodinamika secara absolut dalam memodelkan opasitas plasma \cite{hansen_modeling_2021, favre_merlin_2025}. Namun, paradigma ini secara inheren berujung pada kelumpuhan komputasional ketika diaplikasikan untuk diagnostik analitik. Sebagaimana yang dievaluasi oleh Favre et al. (2025), meskipun model komputasi radiatif murni sangat presisi, penggunaan algoritma pencarian akar non-linear untuk membalikkan matriks RTE secara iteratif memicu hambatan komputasi (computational bottleneck) yang ekstrem. Menjalankan beban integrasi numerik yang kaku pada spektrum eksperimen beresolusi tinggi menuntut biaya waktu yang tidak rasional, sebuah limitasi fundamental yang seketika menggugurkan kelayakan metode analitik ini untuk keperluan diagnostik secara seketika (real-time) \cite{favre_towards_2025}.

% Meskipun substitusi analitik menggunakan algoritma \textit{machine learning} menjanjikan akselerasi komputasi yang signifikan, paradigma statistik murni ini menyimpan cacat epistemologis yang fatal jika diimplementasikan tanpa intervensi fisis. Sebagaimana diekspos secara komprehensif dalam tinjauan Hao et al. \cite{hao_machine_2024}, mayoritas algoritma \textit{data-driven} beroperasi sebagai kotak-hitam (\textit{black-box}) stokastik yang sepenuhnya mencederai kausalitas mekanika kuantum dan memutus rantai hukum termodinamika dasar. Keterputusan antara prinsip fisis dan arsitektur algoritma ini tidak hanya melumpuhkan kemampuan generalisasi model saat berhadapan dengan anomali efek matriks, tetapi juga memicu ketergantungan absolut pada akuisisi dataset empiris berskala masif yang secara eksperimental sangat tidak rasional. Berangkat dari limitasi absolut tersebut, penelitian ini menolak penggunaan \textit{machine learning} sebagai pencocok pola buta, dan memformulasikannya murni sebagai \textit{Surrogate Model} yang ruang pembelajarannya telah dibatasi dan didikte secara deterministik oleh generator fisika pLTE-RTE di fase \textit{offline}.

% Menghadapi kebuntuan tiga arah ini, tujuan utama dari penelitian ini bukanlah merancang algoritma inversi LIBS konvensional, melainkan mengeksekusi resolusi tunggal yang mutlak terhadap Trilema Metodologis tersebut. Perancangan arsitektur \textit{Physics-Informed Surrogate Model} tidak lagi berstatus sebagai opsi analitik, melainkan satu-satunya entitas komputasi yang selamat. Integritas hukum mekanika kuantum sepenuhnya diamankan melalui generator \textit{partial}-LTE (pLTE) dan RTE \textit{offline} yang memproduksi anomali spektral deterministik melalui model dua-zona komprehensif \cite{zaitsev_two-zone_2024}, sementara hambatan komputasi (\textit{computational bottleneck}) dielakkan secara elegan dengan merekam kausalitas tersebut ke dalam proksi \textit{Support Vector Regression} (SVR) \textit{online}. Metodologi hibrida ini memastikan diagnostik ekstraksi termodinamika tereksekusi secara instan, kebal terhadap kesesatan \textit{black-box fallacy}, dan yang terpenting, tidak pernah mendistorsi apalagi mengkhianati fondasi fisika fundamentalnya \cite{favre_merlin_2025}.

% % ==================================================================
\section{Metodologi}
% ============================================================================

Desain metodologis dalam penelitian ini dikembangkan menggunakan Arsitektur Hibrida Fase-Ganda (\textit{dual-phase hybrid architecture}). Berbeda dengan pendekatan analitik konvensional yang memodelkan keseluruhan komponen material secara menyeluruh (\textit{brute-force}), kerangka komputasi ini difokuskan pada pemodelan garis diagnostik utama. Pembangkitan data sintetik (\textit{offline forward modeling}) dilakukan menggunakan pemodelan fisika dasar yang mencakup dimensi transfer radiatif dan keseimbangan pLTE (Subbab \ref{sec:rte} -- \ref{sec:aparatus}). Dataset yang diakumulasikan lalu digunakan untuk melatih prediksi model interpolasi berbasis statistik (\textit{online surrogate modeling}) (Subbab \ref{sec:surrogate}). Pendekatan dua-fase ini dirancang secara sistematis untuk mengatasi lambatnya resolusi numerik iteratif matriks Persamaan Transfer Radiatif (RTE), sekaligus mengurangi ketergantungan empirik pada ekstraksi stokastik \textit{black-box} algoritma \textit{machine learning}.

\subsection{Arsitektur Transfer Radiatif Dua-Zona} \label{sec:rte}

Untuk menangkap tren spasial dalam plasma yang heterogen tanpa menambah berat biaya komputasi di tingkat resolusi matriks penuh matriks resolusi hidrodinamika 3D, emisi makroskopis dideskripsikan ke dalam asumsi distribusi model Dua-Zona yang meliputi area inti emisi dominan (\textit{Core}, $V_1$) serta area marginal penyerapan periferi (\textit{Shell}, $V_p$) \cite{zaitsev_two-zone_2024,hansen_modeling_2021}. Transportasi radiasi sepanjang plasma non-isothermal ini mengikuti Persamaan Transfer Radiatif fundamental secara diferensial komprehensif.
\begin{equation} \label{eq:rte_general}
\frac{dI_\lambda}{ds} = j_\lambda - \kappa_\lambda I_\lambda
\end{equation}

Radiasi zona internal diproyeksikan sebagai intensitas radiasi primer murni $I_{\text{core}}(\lambda)$, yang lantas tereduksi dan melewati parameter absorpsional lokal selubung (parameter ketebalan optik $d_{\text{shell}}$ konstan lokal $j_{\text{shell}}$ dan $\kappa_{\text{shell}}$). Saat integral linear diestimasi pada batasan tepi batas optik tepi.
\begin{equation} \label{eq:rte_integral}
\int_{I_{\text{core}}}^{I_{\text{obs}}} \frac{dI_\lambda}{j_{\text{shell}} - \kappa_{\text{shell}} I_\lambda} = \int_{0}^{d_{\text{shell}}} ds
\end{equation}

Representasi observasi spektrum terdistorsi final dari plasma ($I_{\text{obs}}$) yang merujuk pada analisis transfer analitik ini disajikan presisi menjauhi pendekatan fiktif radiasi gelombang termal \textit{Black Body} \cite{espinosa_analysis_2019}.
\begin{equation} \label{eq:rte_solution}
I_{\text{obs}}(\lambda) = I_{\text{core}}(\lambda) e^{-\tau_\lambda} + \frac{j_{\text{shell}}(\lambda)}{\kappa_{\text{shell}}(\lambda)} \left( 1 - e^{-\tau_\lambda} \right)
\end{equation}
Dimensi optik plasma $\tau_\lambda$ direpresentasikan dari fungsi $\kappa_{\text{shell}} d_{\text{shell}}$, dengan rasio $\frac{j}{\kappa}$ merupakan ekuilibrium parameter tingkat absorpsi yang tidak beranjak sejalan murni dari identitas emisi-absorpsi model asimtotik Kirchhoff  \cite{rodriguez_mixkiprapcal_2021}.

\subsection{Pendekatan Kinetika Partial-LTE (pLTE)}

Metode kinetik konvensional komplit cenderung melacak dinamika perpindahan atomik elektron transien merujuk pada matriks model kolisional-Radiatif (CR) \cite{fujimoto_plasma_2004}.
\begin{equation} \label{eq:master_kinetic}
\frac{dn_i}{dt} = \sum_{j > i} \left( n_j A_{ji} + n_j n_e q_{ji}^{\text{down}} \right) + \sum_{j < i} n_j n_e q_{ji}^{\text{up}} - n_i \sum_{j < i} \left( A_{ij} + n_e q_{ij}^{\text{down}} \right) - n_i \sum_{j > i} n_e q_{ij}^{\text{up}}
\end{equation}
Variabel probabilitas emisivitas didefinisikan lewat Einstein $A_{ji}$, di samping komponen kerapatan dan eksitasi kolisional ($n_e, q^{\text{up}}, q^{\text{down}}$). Namun, literatur komputasi transisi model tingkat eksitasi sekunder mengalami deviasi ketidakpastian (\textit{Atomic Data Desert}) atas parameter mekanika data kolisional atom empirik NIST pada miliaran garis atom \cite{bultel_mcwhirter_2025,favre_merlin_2025}. Terbatasnya data ini menyebabkan matriks $\mathbf{M}$ dalam sistem persamaan laju linear $\frac{d\mathbf{n}}{dt} = \mathbf{M} \cdot \mathbf{n}$ menjadi terlalu jarang (\textit{sparse}), sehingga menghasilkan komputasi distribusi populasi resolusi analitik cacat. Meskipun Fujimoto telah memformulasikan tingkat fase eksitasi matriks CR secara komprehensif \cite{fujimoto_plasma_2004}, ketiadaan data kolisional empiris yang lengkap tersebut memaksa kita untuk mereduksi sistem komputasional ini mutlak kembali menuju batas pLTE.

Sebagai solusi diagnostik kinetik spektrum transien eksperimental, struktur komputasi direstriksi dan disederhanakan murni menjadi struktur termal \textit{Partial}-LTE (pLTE). Pada tekanan densitas ekstrem transien tumbukan awal LIBS, kesetimbangan rasio elektron-ionik dari spesies (\textit{bound-free}) diderivasi secara terkalibrasi melalui Persamaan Saha.
\begin{equation} \label{eq:saha}
\frac{n_{e} n_{II}}{n_{I}} = 2 \frac{Z_{II}(T_e)}{Z_{I}(T_e)} \left( \frac{2\pi m_e k_B T_e}{h^2} \right)^{3/2} \exp\left(-\frac{E_{ion}}{k_B T_e}\right)
\end{equation}
Sedangkan resolusi analitik probablilitas untuk mendistribusi status konfigurasi elektron dari keadaan populasi elektron terikat (\textit{bound excited states}) diatur memakai aproksimasi distribusi eksponensial Boltzmann secara proporsional demi kompensasi kelangkaan korelasi tumbukan empirik sekunder.
\begin{equation} \label{eq:boltzmann_plte}
n_i = \frac{n_{\text{total}} \cdot g_i \exp\left(-E_i / k_B T_e\right)}{Z(T_e)}
\end{equation}
Penerapan pLTE ini memungkinkan terjaganya representasi fisika model serapan optik secara spasial-makroskopis bersamaan dengan penyederhanaan hambatan tingkat penampang tabrakan mikro atomikal internal komputasi secara mandiri \cite{cristoforetti_local_2010}.

\subsection{Asimilasi Parameter Garis Diagnostik Terseleksi}

Distribusi emisi $j$ dan serapan radiatif $\kappa$ yang tertuang sebagai rasio makroskopik (Persamaan \ref{eq:rte_solution}) dikonstruksi secara mikroskopis oleh komponen probabilitas termal non-LTE.
\begin{equation} \label{eq:emission_coeff}
j_\lambda = \sum_{k \to i} \frac{hc}{4\pi\lambda} A_{ki} n_k \phi_V(\lambda)
\end{equation}
Di samping itu peluruhan optik korelasi radiatif akibat foton yang menstimulasi laju penurunan pita eksitasi (\textit{detailed balance}) mempertegas besaran nilai konvolusi fungsi absorpsi murni \cite{espinosa_analysis_2019}.
\begin{equation} \label{eq:absorption_coeff}
\kappa_\lambda = \sum_{i \to k} \frac{\lambda^4}{8\pi c} \left(\frac{g_k}{g_i}\right) A_{ki} n_i \phi_V(\lambda) \left[ 1 - \frac{n_k g_i}{n_i g_k} \right]
\end{equation}

Pemodelan spektral komputasi dalam studi ini difokuskan secara independen pada identifikasi garis parameter resonansi yang telah tervalidasi secara eksperimental. Metode komputasi teoretis komprehensif pada dasarnya cenderung mengekstrapolasi data dan menyatukan seluruh kemungkinan lintasan resolusi atomik sekunder secara polinomial \cite{favre_merlin_2025}, yang berisiko menciptakan deviasi margin rentan akibat profil simulasi yang tidak reliabel. Oleh karenanya, struktur simulasi memprioritaskan "Garis Diagnostik Utama" (contohnya doublet resonansi elemen Ca II dan Si I) di mana setiap landasan mekanika kuantum emisiloginya dapat dirujuk langsung pada literatur spektroskopi definitif empiris \cite{griem_spectral_1974}. Skema penetapan diagnostik \textit{benchmark} resolusi spesifik elemen bertujuan mengisolasi data set generatif (\textit{database}) murni dalam kaidah batas matriks akurasi fisis agar algoritma tipe \textit{Surrogate} di babakan fase lanjutan dapat tereksekusi tanpa interferensi korelasi stokastik parameter prediksi minor \cite{tang_review_2024, manelski_libs_2024}. 

\subsection{Konvolusi Profil Voigt dan Resolusi Aparatus} \label{sec:aparatus}

Distribusi profil penyebaran energi termodinamika fisis pada pita frekuensi fungsi spektroskopi linear dirangkum di dalam konvolusi perpaduan kurva fungsi matematis $\phi_V(\lambda)$. Komposisi konvolusi dihitung berdasar efek pelebaran Doppler terkait profil fungsi asimetris Gaussian gerak elektron termal.
\begin{equation} \label{eq:broadening_doppler}
\Delta\lambda_D = \lambda_{ki} \sqrt{8 \ln 2 \frac{k_B T_i}{m_k c^2}}
\end{equation}
Serta kombinasi interaksi elektrostatis resolusi turbulensi plasma tekanan LIBS (\textit{Stark Broadening}) dominan yang mencirikan profil Lorentzian pada batasan analit densitas elektron $n_e \ge 10^{16}$ cm$^{-3}$ sejalan dengan rumusan pelebaran mekanika kuantum empiris oleh Griem \cite{griem_spectral_1974}.
\begin{equation} \label{eq:broadening_stark_simp}
\Delta\lambda_S \approx 2 \omega_{ki}(T) \frac{n_e}{n_{e,r}}
\end{equation}
Perpaduan fungsi densitas suhu dan parameter \textit{Stark} diekstraksi ke dalam konvolusi Voigt $\phi_V(\lambda)$. Struktur fungsi Voigt bersama formula RTE Dua-Zona (Persamaan \ref{eq:rte_solution}) merumuskan deformasi intensitas puncak spektral akibat distorsi medium optik tebal plasma (\textit{self-reversal}). 

Lebih lanjut, eksperimen pendeteksian optik nyata turut memudarkan (\textit{smearing effect}) presisi kecekungan pelebaran absorpsi fraunhofer optik fisis akibat interferensi limitasi fungsi difraksi aparatus mikrometer spektrometer (\textit{slit} difraksi analitik) \cite{zaitsev_two-zone_2024}. Bentuk empirik akhir dari konvolusi analitik resolusi aparatus tersebut memodifikasi sebaran simulasi ($I_{obs}$) dengan distribusi piksel Gaussian ($G_{inst}$) rata-rata alat ($\Delta\lambda_{inst} \approx 0.1$ nm).
\begin{equation} \label{eq:instrumental}
I_{final}(\lambda) = I_{obs}(\lambda) \otimes G_{inst}(\lambda, \Delta\lambda_{inst})
\end{equation}
Prosedur resolusi aparatus instrumen akhir menjembatani perbedaan parameter teoretik agar secara ekuivalen terdistribusi sejajar menyamai gradien profil kecekungan spektrum terekam \cite{favre_merlin_2025}.

\subsection{Pelatihan Surrogate Model Berbasis SVR} \label{sec:surrogate}

Fase ekstraktif algoritma komputasional arsitektur diagnostik plasma dilaksanakan sebagai skema pembalikan resolusi parameter linier di tahap akhir.

Memasuki fase komputasi murni (\textit{offline}), generator sistem simulasi memanifestasi formulasi ekuilibrium persamaan RTE, konvolusi pLTE, pelebaran difraksi optik (\textit{Voigt-aparatus}), menjadi representasi sistem periodik. Parameter transien plasma kemudian dibangun dengan eksekusi literasi $T_e$ mulai kisaran nilai [4000, 20000]\,K terhadap $n_e \in [10^{14}, 10^{18}]$\,cm$^{-3}$ difokuskan memvalidasi set garis divalen Ca II atom komposit campuran Si I. Data simulatif termodinamik ini menghasilkan domain ematur pra-kondisi terkurasi fisis deterministik pLTE. 

Kumpulan emisi transien tersebut diadaptasikan guna melatih matriks regresi asimtotik Support Vector Regression (SVR) secara spesifik (fase \textit{online}) \cite{wu_novel_2025, lin_enhancing_2026}. Model pembelajaran regresi diturunkan sebagai proksi interpolasi instan linear (\textit{surrogate}) atas nilai pembalikan identitas emisi yang sebetulnya dikonstruksi berbasis fisika fundamental. Pendekatan percontohan termodinamika memitigasi kompleksitas resolusi perhitungan deterministik iteratif dalam penyelesaian akar inversi RTE menjadi evaluasi parametrik algoritma matematis prediktif analitis instan (sekala diagnostik \textit{real-time}). Realisasi proklamasi kecepatan asimtotik deterministik termodinamik SVR menjanjikan keakuratan pengukuran analit observasional profil garis padat tanpa membedakan nilai empiris laboratorium optik tebal plasma.

% ============================================================================================
\bibliographystyle{elsarticle-num}
\bibliography{reff}

\section*{References}



\end{document}


